%% ============================================================
%% 分片 19：原文第 7 节「Fluxes and $S$-Duality」（原文 glshort.tex 起于第 5685 行）
%% 由 tools/harvmac2tex.py 从 arXiv 源码（hep-th/0604151v3）自动生成。
%% 只译英文散文；公式、编号、\label、\cite 一律不改。
%% ============================================================

\section{Fluxes and $S$-Duality}

\label{sec:centertop}

In this section, we will examine phenomena involving the center of
the gauge group that are relevant to understanding $S$-duality of
Hitchin moduli space and the geometric Langlands program.

We begin section 第~\ref{sec:reviewthis}~节\ with a review of discrete electric
and magnetic fluxes and the action of $S$-duality on them. More
detail can be found in \cite{ref45} , section 3, as well as in
\cite{ref26}. For discrete electric and magnetic flux in gauge theory,
the original reference is \cite{ref101}. We also explain some mathematical notions such as the
concept of the universal bundle. Then we develop our applications
in section 第~\ref{sec:comptwop}~节.

\subsection{Review}

\label{sec:reviewthis}

In this section, $G$ denotes a compact semi-simple Lie group;
$\bar G$ is its universal cover; and $G_{\rm {ad}}$ is the
corresponding adjoint group.  The center of $\bar G$, which is
also the fundamental group of $G_{\rm {ad}}$, is a finite abelian
group $\CZ$.  For simple $G$, $\CZ$ is a cyclic group in all cases
except for $\bar G=Spin(4n)$, as is summarized in the table.

One happy fact that we notice from the table is that some of the
subtleties with $S$-duality for non-simply-laced groups are
irrelevant, since $G_2$ and $F_4$ have trivial centers.

\par\medskip

\begin{table}[htbp]
\centering
\begin{tabular}{ll}
\toprule
Group & Center \\
\midrule
$SU(N)$ & $\Bbb{Z}_N$ \\
$Spin(2n+1)$ & $\Bbb{Z}_2$ \\
$Spin(4n+2)$ & $\Bbb{Z}_4$ \\
$Spin(4n)$ & $\Bbb{Z}_2\times\Bbb{Z}_2$ \\
$Sp(n)$ & $\Bbb{Z}_2$ \\
$E_6$ & $\Bbb{Z}_3$ \\
$E_7$ & $\Bbb{Z}_2$ \\
\bottomrule
\end{tabular}
\caption{Simple and simply-connected Lie groups with non-trivial centers. $F_4$, $G_2$, and $E_8$, which are omitted, have trivial center.（原文 Table~3；原文此处误作 $F_2$、$G_4$）}
\label{tab:centers}
\end{table}

An important property  is that the group ${\cal Z}$ is naturally
selfdual.  The dual of an abelian group $B$ is $B^\vee={\rm
Hom}(B,U(1))$.  A finite abelian group is always isomorphic to its
own dual, but not naturally.  The center $\CZ$ of  $\bar G$ always
has, however, a natural selfduality. This duality is best
expressed as a homomorphism % 式 (7.1)
\setcounter{equation}{0}
\begin{equation}\label{eq:plook}
\Upsilon:\CZ\times \CZ\to
U(1),
\end{equation}
 which is symmetric, $\Upsilon(a,b)=\Upsilon(b,a)$, and such
that any homomorphism of $\CZ$ to $U(1)$ is $a\to \Upsilon(a,b)$
for a unique $b$.  Such a $\Upsilon$ is called a ``perfect
pairing.'' $\Upsilon$ is constructed as follows.  If $G$ is
simply-laced, then $\CZ=\Lambda^\vee/\Lambda$, where $\Lambda$ is
the root lattice and $\Lambda^\vee$ is the weight lattice of $\bar
G$.  The pairing $\Upsilon$ is then defined as
% 式 (7.2)
\setcounter{equation}{1}
\begin{equation}\label{eq:rofo}
\Upsilon(a,b)=\exp(2\pi i\langle a,b\rangle),
\end{equation}
 where
$\langle~,~\rangle$ is the usual quadratic form on $\Lambda^\vee$
(which is even and integral when restricted to $\Lambda$ and takes
rational values on $\Lambda^\vee$). If $G$ is not simply-laced,
its center is trivial or $\Bbb{Z}_2$, and so admits precisely one
selfduality, which we call $\Upsilon$.

 A $\Gad$-bundle $E$ on $M$ has a characteristic class $\xi(E)$
that takes values in $H^2(M,\CZ)$. For example, if $G_{{\rm
ad}}=SO(3)$, $\xi(E)$ is the second Stieffel-Whitney class
$w_2(E)$.  One can define a partition function $Z_\xi$ for every
choice of $\xi$ by restricting the path integral to bundles with
given $\xi$.  The $Z_\xi$ transform in a unitary representation of
the $S$-duality group \cite{ref45,ref26}, but this is
something that we will not explain here.

Instead, we concentrate on the Hamiltonian description, and
specialize to $M=S^1\times W$ for a three-manifold $W$. In this
case, we have $H^2(M,\CZ)=H^1(W,\CZ)\oplus H^2(W,\CZ)$, so
$\xi(E)$ can be decomposed as $\xi={\bf a}\oplus {\bf m}$, with
${\bf a}\in H^1(W,\CZ)$, ${\bf m}\in H^2(W,\CZ)$.

Consider the pairing % 式 (7.3)
\setcounter{equation}{2}
\begin{equation}\label{eq:junko}
\hat \Upsilon:H^2(M,\CZ)\times
H^2(M,\CZ)\to H^4(M,U(1))=U(1)
\end{equation}
 obtained by composing the cup
product $H^2(M,\CZ)\times H^2(M,\CZ)\to H^4(M,\CZ\times\CZ)$ with
the map $\Upsilon:\CZ\times\CZ\to U(1)$. According to Poincar\'e
duality, $\hat\Upsilon$ is a perfect pairing, making $H^2(M,\CZ)$
a selfdual abelian group. More generally, for a closed oriented
manifold $Y$ of dimension $n$, Poincar\'e duality together with
selfduality of $\CZ$ gives perfect pairings
% 式 (7.4)
\setcounter{equation}{3}
\begin{equation}\label{eq:unko}
H^d(Y,\CZ)\times H^{n-d}(Y,\CZ)\to U(1), ~~0\leq d\leq
n.
\end{equation}
 In particular, we have the perfect pairing
% 式 (7.5)
\setcounter{equation}{4}
\begin{equation}\label{eq:normo}
\bar\Upsilon:H^1(W,\CZ)\times H^2(W,\CZ)\to U(1).
\end{equation}
 So
these are dual abelian groups, and in particular
% 式 (7.6)
\setcounter{equation}{5}
\begin{equation}\label{eq:junok}
H^2(W,\CZ)={\rm Hom}(H^1(W,\CZ),U(1)).
\end{equation}

Let us now recall how $S$-duality is implemented in a Hamiltonian
framework. The partition function of the gauge theory with ${\bf
m}$ specified has a natural Hilbert space interpretation, because
${\bf m}$ is part of the data at an initial time and is
independent of time. However, because ${\bf a}$ cannot be
expressed in terms of the data at a fixed time, the partition
function with fixed ${\bf a}$ cannot be given a Hamiltonian
interpretation.  Instead  one must introduce a character ${\bf e}$
of the finite abelian group $H^1(M,\CZ)$, that is, a homomorphism
% 式 (7.7)
\setcounter{equation}{6}
\begin{equation}\label{eq:kilno}
{\bf e}:H^1(W,\CZ)\to U(1).
\end{equation}
 For each choice of ${\bf
e}$ and ${\bf m}$, the sum $\sum_{{\bf a}}{\bf e}({\bf a})Z_{{\bf
a},{\bf m}}$ can be interpreted in terms of a trace in a Hilbert
space ${\cal H}_{{\bf e},{\bf m}}$. This sum has a Hilbert space
interpretation because summing over ${\bf a}$ with the weight
factor ${\bf e}({\bf a})$ is compatible with cutting and pasting
(or what  in quantum field theory is usually called cluster
decomposition) in the $S^1$ direction. ${\bf e}$ and ${\bf m}$ are
called respectively the discrete electric and magnetic flux
\cite{ref101}. Because of \eqref{eq:junok}, $\bf e$ can be alternatively viewed as
an element of $H^2(W,\CZ)$, just like $\bf m$.

So we get a Hilbert space ${\cal H}_{{\bf m},{\bf e}}$ for each
choice of % [WARN] 公式 \eqn\eqref{eq:ytro} 无编号或重复
\begin{equation*}
\tag*{(7.8)}
\begin{aligned}{\bf e}&\in {\bf E}={\rm
Hom}(H^1(W,\CZ), U(1))\\
                             {\bf m}& \in {\bf M}=H^2(W,\CZ).\\\end{aligned}
\end{equation*}

Happily, the finite groups ${\bf E}$ and ${\bf M}$ are isomorphic
according to \eqref{eq:junok}. This makes it possible to have $S$-duality
symmetry. The duality transformation $S$ exchanges ${\bf E}$ and
${\bf M}$ and maps $({\bf e},{\bf m})\to (-{\bf m},{\bf e})$:
% 式 (7.9)
\setcounter{equation}{8}
\begin{equation}\label{eq:coco}
S:{\cal H}_{{\bf e},{\bf m}}\to {\cal H}_{-{\bf m},{\bf
e}}.
\end{equation}
 It also, of course, maps $\tau\to -1/\tau$.  (Eqn. \eqref{eq:coco}\ is
a discrete analog of \eqref{eq:xilj}, with $(\vec{\eurm m},\vec{\eurm n})$
replaced by their discrete counterparts $({\bf m},{\bf e})$.)

While it is possible to define all the Hilbert spaces ${\cal
H}_{{\bf e},{\bf m}}$, a given gauge theory construction may not
use all of them.  If we do gauge theory on $M=S^1\times W$ with
simply-connected gauge group $\bar G$, then we must set to zero
the characteristic classes ${\bf a}$ and ${\bf m}$ that enter in
the partition function $Z_{{\bf a},{\bf m}}(M)$. For the Hilbert
space, this means that ${\bf e}$, which arises by a Fourier
transform with respect to ${\bf a}$, is arbitrary, while ${\bf
m}=0$.  Hence % 式 (7.10)
\setcounter{equation}{9}
\begin{equation}\label{eq:ytoto}
{\cal H}(\bar G,W)=\oplus_{{\bf e}}{\cal
H}_{{\bf e},0}.
\end{equation}
 Taking into account the assumed transformation
law for ${\bf e},{\bf m}$ under $S$-duality, we see that the
Hilbert space of the dual theory must be % 式 (7.11)
\setcounter{equation}{10}
\begin{equation}\label{eq:momjo}
\oplus_{{\bf
m}}{\cal H}_{0,{\bf m}}.
\end{equation}
 How is this Hilbert space related to a
path integral on $M=S^1\times W$? Setting $\bf e$ to zero means
summing over $\bf a$. So in terms of the variables ${\bf a},{\bf
m}$, \eqref{eq:momjo}\ means one must sum over all ${\bf a},{\bf m}$ with
equal weight. In other words, in the dual theory we sum over
$^L\Gad$-bundles with all possible $\xi$. Thus the dual theory has
gauge group $^L\Gad$, as expected.

\par\medskip{\it The Universal Bundle}

We need another piece of background, which is the concept of the
universal bundle.

Consider first the moduli space $\M(G,C)$ of $G$-bundles over a
Riemann surface $C$.  For every point $p\in \M(G,C)$, there is a
corresponding $G$-bundle $E_p\to C$ parametrized by $p$; it is
determined up to isomorphism.  A universal bundle, if it exists,
is a $G$-bundle ${\EUE} \to \M\times C$ such that for any $p\in
\M$, ${\EUE}$ restricted to $p\times C$ is isomorphic to $E_p$.
Gauge theory gives \cite{ref91}\ a very direct method to analyze the
universal bundle, though we will not explain this here.

Naively, we just define ${\EUE}$ so that its restriction to
$p\times C$ is ``the'' bundle over $C$ that is determined by $p$.
The reason that a problem arises is that, to begin with, $p$ only
determines an isomorphism class of $G$-bundle over $C$, not an
actual bundle. This causes no difficulty locally\footnote{This statement holds
near a smooth point in $\M$.  A more precise
analysis than we will give shows that the class
$\zeta(\EUE_{\rm ad})$ that obstructs the universal bundle can have a local contribution
at singularities.};
we just pick a
bundle in the right isomorphism class at a particular $p$ and then
deform it in a small neighborhood. So for each small open set
$U_i\subset \M$, we can pick a universal bundle ${\EUE}_i\to
U_i\times C$, and moreover ${\EUE}_i$ is unique up to isomorphism.

The next step is to glue together the ${\EUE}_i$ over $ (U_i\cap
U_j)\times C$.  The bundles ${\EUE}_i$ and ${\EUE}_j$ are
isomorphic over $ (U_i\cap U_j)\times C$, so we pick an
isomorphism $\Theta_{ij}:{\EUE}_i\cong {\EUE}_j$ over this
intersection, with $\Theta_{ji}=\Theta_{ij}^{-1}$.

Now if on triple intersections $(U_i\cap U_j\cap U_k)\times C$,
the composition $\Theta_{ki}\Theta_{jk}\Theta_{ij}$ is equal to 1,
then we can consistently glue together the ${\EUE}_i$ to get the
desired universal bundle ${\EUE}$.

If the gauge group is the adjoint group $G_{{\rm ad}}$, there is
no problem in the gluing. A generic $G_{{\rm ad}}$ bundle has no
automorphisms (exceptions occur at singularities of $\M$). So
$\Theta_{ki}\Theta_{jk}\Theta_{ij}$, just because it is an
isomorphism of ${\EUE}_i$ (restricted to $(U_i\cap U_j\cap
U_k)\times C$) to itself, is the identity at a generic point of
the triple intersection and hence everywhere.

Thus, the universal   $G_{{\rm ad}}$-bundle ${\EUE}_{\rm {ad}}$
does exist. Now let us repeat this discussion, taking the gauge
group to be the simply-connected cover $\bar G$ (or more
generally, any cover of $G_{\rm{ad}}$ with a non-trivial center).
We start in the same way. For each small open set $U_i\in \M$, we
pick a universal bundle $\bar {\EUE}_i$ over $U_i\times C$, and
try to glue to make a universal $\bar G$-bundle $\bar {\EUE}$.

The reason that the result is different is that a generic $\bar
G$-bundle does have a non-trivial group of automorphisms; the
center $\CZ$ is a group of automorphisms of any $\bar G$-bundle
$E$.

Hence, in the above argument, the composite map
$\Theta_{ki}\Theta_{jk}\Theta_{ij}$ over a triple intersection is
not necessarily the identity; instead % 式 (7.12)
\setcounter{equation}{11}
\begin{equation}\label{eq:yrf}
\Theta_{ki}\Theta_{jk}\Theta_{ij}=f_{ijk},
\end{equation}
 with
$f_{ijk}\in\CZ$. The $f_{ijk}$ combine to a two-cycle defining an
element $\zeta\in H^2(\M,\CZ)$. The element $\zeta$ is known to be
non-trivial; it is the obstruction to the existence of the
universal $\bar G$-bundle $\bar{\EUE}$.

There is another way to explain the existence of $\zeta$.  Let
$\M(G_{\rm{ad}},C)_0$
  be the component of $\M(G_{\rm{ad}},C)$ that
parametrizes topologically trivial bundles -- the ones that when
restricted to $p\times C$, for $p\in \M(G_{\rm{ad}},C)$, can be
lifted to $\bar G$-bundles.    Consider the universal bundle
${\EUE}_{\rm{ad}}$ over   $\M(G_{\rm{ad}},C)_0\times C$.  If we
could lift ${\EUE}$ to a $\bar G$-bundle $\bar {\EUE}$, this
(after being pulled back to $\M(\bar G,C)\times C$, which is a
finite cover of $\M(G_{\rm {ad}},C)_0\times C$) would be the
universal $\bar G$-bundle. In general, the obstruction to lifting
a $G_{\rm ad}$-bundle ${\EUE}_{\rm{ad}}$ to a $\bar G$-bundle
$\bar {\EUE}$ is the characteristic class $\xi({\EUE}_{\rm{ad}})$.
So the obstruction to existence of the universal $\bar G$-bundle
is $\zeta=\xi({\EUE}_{\rm {ad}})$.

There is a similar story for Higgs bundles.  The universal Higgs
bundle is a pair $( {\EUE},\hat \varphi)$,
 with $ {\EUE} $ being a $G$-bundle over
$ \MH(G,C)\times C$, and $\hat\varphi\in H^0( \MH(G,C)\times
C,{\rm {ad}}({\EUE})\otimes K_C)$, obeying the following
condition. For each $p\in \MH(G,C)$, the restriction of $(
{\EUE},\hat\varphi)$ to $p\times C$ should be isomorphic to the
Hitchin pair $(E_p,\varphi_p)$ parametrized by $p$.

The arguments that we have already given can be carried over with
no essential change to show that for gauge group $G_{\rm{ad}}$,
the universal Higgs bundle $( {\EUE}_{\rm {ad}},\hat\varphi)$ does
exist. But for the simply-connected gauge group $\bar G$ (or any
nontrivial cover of $G_{\rm {ad}}$), the universal Higgs bundle
does not exist. It is obstructed by the fact that a generic
Hitchin pair $(E,\varphi)$, with $E$  a $\bar G$-bundle, has the
group $\CZ$ of automorphisms. The obstruction is just the
obstruction to lifting the ``bundle'' part of the universal
$G_{{\rm ad}}$ Higgs bundle $({\EUE}_{\rm{ad}},\hat\varphi)$ to a
$\bar G$-bundle. It is therefore
% 式 (7.13)
\setcounter{equation}{12}
\begin{equation}\label{eq:plooo}
\zeta=\xi({\EUE}_{\rm{ad}}).
\end{equation}

\par\medskip{\it The Universal Bundle As A Twisted Vector
Bundle}

Although for $G$ not of adjoint type, the universal $G$-bundle
does not exist as a vector bundle or as a principal bundle, it
does exist as a twisted vector bundle.  We will describe this
concept in a very pedestrian way.

 One way to construct an ordinary vector
bundle $V$ of rank $N$ over a manifold $X$ is to cover $X$ with
small open sets $U_i$, on each of which we pick a rank $N$ trivial
bundle $V_i\cong U_i\times \Bbb{C}^N$. Then we glue $V_i$ to $V_j$
on the intersection $U_i\cap U_j$ via a gluing map $v_{ij}:V_i\cap
V_j\to U(N)$, with $v_{ji}=v_{ij}^{-1}$. If on triple
intersections we have % 式 (7.14)
\setcounter{equation}{13}
\begin{equation}\label{eq:nofto}
v_{ki}v_{jk}v_{ij}=1,
\end{equation}
 then the
$V_i$ can be glued together consistently to make a rank $N$ vector
bundle $V\to X$.

Now suppose that we are given $b\in H^2(X,U(1))$.  Then $b$ can be
represented explicitly by a $U(1)$-valued cocycle $b_{ijk}$
defined on triple intersections $U_i\cap U_j\cap U_k$. A twisted
vector bundle is defined by the same sort of data $v_{ij}:V_i\cap
V_j\to U(N)$.  But now, instead of \eqref{eq:nofto}, we ask for
% 式 (7.15)
\setcounter{equation}{14}
\begin{equation}\label{eq:hofto}
v_{ki}v_{jk}v_{ij}=b_{ijk}.
\end{equation}
 Thus, a $b$-twisted vector
bundle $V$ is not a vector bundle in the usual sense.  However,
the associated adjoint bundle ${\rm ad}(V)$ is an ordinary vector
bundle, since the phase $b_{ijk}$ in \eqref{eq:hofto}\ disappears if we pass
to the adjoint representation.

Now let us return to the problem of constructing a universal
bundle.  In this paper, we are generally a little imprecise about
whether by a $G$-bundle, we mean a principal $G$-bundle, whose
fiber is a copy of $G$, or an associated vector bundle in some
faithful  representation of $G$.  Principal bundles make possible
a uniform analysis good for any $G$, but for a group like $U(N)$
that has a convenient faithful representation (the $N$-dimensional
representation) it is useful to think in terms of vector bundles.

In discussing the universal bundle, it is helpful to be more
precise.  We interpret the transition functions $\Theta_{ij}$ of
our discussion above as $\bar G$-valued functions, with no
particular choice of representation.  Their projection to $G_{ad}$
gives us transition functions for a universal principal
$G_{ad}$-bundle ${\EUE}_{\rm ad}$, but we cannot lift this to a
principal $\bar G$-bundle, because of the relation
% 式 (7.16)
\setcounter{equation}{15}
\begin{equation}\label{eq:flutter}
\Theta_{ki}\Theta_{jk}\Theta_{ij}=f_{ijk},
\end{equation}
 where the
$f_{ijk}$ define the class $\zeta=\xi({\EUE}_{ad})\in
H^2(\MH(G_{\rm ad},C),\CZ)$. Now let us pick an irreducible
representation $\varrho$ of $\bar G$.  We write
$\Theta_{ij}^\varrho$ for the transition functions $\Theta_{ij}$
evaluated in the representation $\varrho$. Likewise we write
$\varrho(f)$ for $f$ evaluated in the representation $\varrho$. In
the irreducible representation $\varrho,$ the center of $\bar G$
acts by scalar multiplication, so $\varrho(f)$ takes values in
$U(1)$.  We have
% 式 (7.17)
\setcounter{equation}{16}
\begin{equation}\label{eq:glutter}
\Theta^\varrho_{ki}\Theta^\varrho_{jk}\Theta^\varrho_{ij}=\varrho(f_{ijk}).
\end{equation}

The quantities $\varrho(f_{ijk})$ are a cocycle defining the
element $\varrho(\zeta)\in H^2(\MH(G_{\rm ad},C),U(1))$.

If $\varrho(f)=1$, the objects $\Theta^\varrho$ are transition
functions that define a vector bundle ${\EUE}_\varrho\to \MH(\bar
G,C)\times C$ that we may call the universal bundle in the representation
$\varrho$.  In general this is not the case.  However, comparing
eqns. \eqref{eq:flutter}\ and \eqref{eq:hofto}, we see that while ${\EUE}_\varrho$ may
not exist as an ordinary vector bundle, it does always exist as a
twisted vector bundle, twisted by $\varrho(\zeta)\in
H^2(\MH(G_{\rm ad},C),U(1))$.
